\documentclass[../main.tex]{subfiles}
\usepackage{amsmath}
\usepackage{hyperref}
\usepackage{enumitem}

\begin{document}



\section{Computer Vision}

\subsection{Transferability}
\begin{itemize}

    \item \textbf{Cross-architectural transferability}: Investigating adaptive noise patterns exploiting shared backbone vulnerabilities (e.g., YOLO$\to$Faster R-CNN) using randomized scaling/offsets \cite{cross_model_transfer}.
\end{itemize}

\subsection{Object Detectors and Physical-World Robustness}
\begin{itemize}
    \item \textbf{Bounding-box-constrained perturbations}: Developing attack methods that restrict perturbations to object bounding boxes (e.g., using RCEN or logit loss optimization), minimizing background noise while maintaining high attack success rates ($>$80\% transferability) \cite{bounding_box_attacks}.
    \item \textbf{Temporal attack consistency}: Extending adversarial attacks to video segmentation/tracking systems by integrating hard-region discovery with temporal perturbation constraints \cite{video_attacks}.
    \item \textbf{Physical-world validation}: Testing 3D-printed adversarial objects under occlusion/lighting variations using multiview rendering pipelines \cite{3d_physical_attacks}.
\end{itemize}

\subsection{Defense-Aware Attack Generation}
\begin{itemize}
    \item \textbf{Ensemble defense evasion}: Designing attacks targeting shared vulnerabilities across model architectures (e.g., ResNet, ViT) to bypass adversarial training ensembles \cite{ensemble_evasion}.
    \item \textbf{Geometric-invariant perturbations}: Integrating rotation/scaling invariance into attack generation to counter preprocessing defenses \cite{geometric_invariance}.
\end{itemize}

\section{Language Models}

\subsection{LLM-Powered Adversarial Engines}
\begin{itemize}
    \item \textbf{Multi-modal attack synthesis}: Leveraging LLMs to generate semantically coherent adversarial examples across text, audio, and code domains, exploiting prompt injection vulnerabilities \cite{multimodal_attacks}.
    \item \textbf{Safety-alignment evasion}: Developing targeted attacks overriding ethical constraints in instruction-tuned LLMs using universal adversarial images \cite{safety_alignment_evasion}.
\end{itemize}

\subsection{Cross-Modal and Physical Attacks}
\begin{itemize}
    \item \textbf{Vision-language universal triggers}: Engineering unified adversarial triggers effective across ASR, sentiment analysis, and machine translation \cite{over_air_asr}.
    \item \textbf{Physical ASR robustness}: Designing over-air audio attacks using expectation-over-transformation (EoT) methods simulating real-world distortions \cite{eot_methods}.
\end{itemize}



These directions address critical gaps in attack transferability, real-world applicability, and defense interoperability, advancing toward secure, attack-resistant AI systems in critical applications.

\ifsubfiles
    \bibliographystyle{unsrt}
    \bibliography{references}
\else
    % Empty block for main document compatibility
\fi

\end{document}